\documentclass[aspectratio=169]{beamer}

\usepackage[utf8]{inputenc}
\usepackage[english]{babel}
\usepackage{booktabs}
\usepackage{framed}
\usepackage{graphicx}
\usepackage{tikz}
\usetikzlibrary{arrows.meta,positioning}

\usepackage{empheq}
\usepackage[most]{tcolorbox}

\newtcbox{\mymath}[1][]{%
    nobeforeafter, math upper, tcbox raise base,
    enhanced, colframe=blue!30!black,
    colback=blue!30, boxrule=1pt,
    #1}

\graphicspath{{logo/}{../logo/}}

\definecolor{ForestGreen}{RGB}{22,100,70}
\definecolor{SoftGreen}{RGB}{229,244,235}
\definecolor{SoftBlue}{RGB}{226,239,255}
\definecolor{SoftOrange}{RGB}{255,241,220}
\definecolor{SoftPurple}{RGB}{239,231,255}

\mode<presentation>{
  \usetheme{Ilmenau}
  \usecolortheme{spruce}
  \usefonttheme{professionalfonts}
  \setbeamertemplate{navigation symbols}{}
  \setbeamertemplate{headline}{}
  \setbeamertemplate{caption}[numbered]
  \setbeamercolor{institute in head/foot}{fg=ForestGreen}
  \setbeamercolor{metric box}{bg=ForestGreen,fg=white}
}

\tikzset{
  concept/.style={draw=ForestGreen, rounded corners, fill=SoftGreen,
    very thick, minimum height=1.15cm, text width=3.6cm, align=center},
  application/.style={draw=blue!60!black, rounded corners, fill=SoftBlue,
    thick, minimum height=0.95cm, text width=3.8cm, align=center},
  stage/.style={draw=ForestGreen, rounded corners, fill=SoftGreen,
    thick, minimum height=1.15cm, text width=2.05cm, align=center,
    font=\small},
  arrow/.style={-Latex, very thick, ForestGreen}
}

\title[TOKT11105 - Lecture 1]{Data Structures and Algorithms in Python\\Lecture 1: Introduction}
\author{Faculty of Data Science and Artificial Intelligence}
\institute{College of Technology - National Economics University}
\date{}

\AtBeginSection[]{
  \begin{frame}[plain]
    \begin{center}
      \vfill
      {\color{ForestGreen}\large\bfseries SECTION \insertsectionnumber}

      \vspace{0.8em}
      {\color{ForestGreen}\rule{0.68\textwidth}{1.2pt}}

      \vspace{1.1em}
      {\usebeamerfont{section title}\color{ForestGreen}\Huge\bfseries
        \insertsectionhead}

      \vspace{1.1em}
      {\color{ForestGreen}\rule{0.68\textwidth}{1.2pt}}
      \vfill
    \end{center}
  \end{frame}
}

\begin{document}

\begin{frame}[plain]
  \setbeamercolor{block body example}{bg=ForestGreen,fg=white}
  \begin{center}
    \begin{beamercolorbox}[wd=\textwidth,sep=1pt,center,rounded=true,shadow=true]{block body example}
      \includegraphics[width=0.27\textwidth]{logo_slide.png}
    \end{beamercolorbox}
  \end{center}
  \vspace{-0.25em}
  \titlepage
\end{frame}

\logo{\includegraphics[width=0.8cm]{logo_fda.png}}
\setbeamertemplate{footline}[miniframes theme]

\begin{frame}{}
  \begin{columns}[T]
    \begin{column}{0.49\textwidth}
      \begin{figure}
        \centering
        \includegraphics[width=0.8\linewidth]{figures/AI_code_1.jpg}
      \end{figure}
    \end{column}
    \begin{column}{0.49\textwidth}
      \begin{figure}
        \centering
        \includegraphics[width=0.78\linewidth]{figures/AI_code_2.jpg}
      \end{figure}
    \end{column}
  \end{columns}
\end{frame}

\begin{frame}{Some Facts}
  AI can help write code, but ...
  \begin{enumerate}
    \item understanding why and how fast it runs is still on us.
    \item it cannot replace clear reasoning, pattern recognition, and decision-making.
  \end{enumerate}

  \vspace{0.5cm}

  \fcolorbox{black}{yellow!20}{%
    \begin{minipage}{\dimexpr\linewidth-2\fboxsep-2\fboxrule}
      {\bf Your value lies not in writing code, but in thinking critically---analyzing problems, designing solutions, and guiding AI to produce the right code.}
    \end{minipage}%
  }

  \vspace{0.5cm}

  \noindent ... {\em and that's why you should learn}

  \vspace{0.3cm}

  \begin{empheq}[box=\tcbhighmath]{equation*}
    {\Large\text{\bf Data Structures and Algorithms}}
  \end{empheq}
\end{frame}

\begin{frame}{Today: the big picture first}
  \begin{enumerate}
    \item What is a data structure?
    \item What is an algorithm?
    \item Why are they closely related? Why are they important?
  \end{enumerate}
\end{frame}

\begin{frame}{A short definition}
  \begin{columns}[T]
    \begin{column}{0.5\textwidth}
      \begin{block}{Data structure}
        A way to \textbf{organize and store data} so it can be accessed and modified effectively.
      \end{block}
    \end{column}
    \begin{column}{0.5\textwidth}
      \begin{alertblock}{Algorithm}
        A finite, unambiguous sequence of \textbf{steps for solving a problem}.
      \end{alertblock}
    \end{column}
  \end{columns}

  \vspace{1cm}

  \begin{empheq}[box=\tcbhighmath]{equation*}
    {\LARGE\text{{\color{red}Data Structures} + {\color{blue}Algorithms} = Program}}
  \end{empheq}
\end{frame}

\begin{frame}{Quick check: data structures and algorithms}
  \small

  \begin{block}{Question 1}
    Which statement best describes a \textbf{data structure}?
    \begin{enumerate}
      \item A sequence of steps for solving a problem.
      \item A way to organize and store data for efficient operations.
      \item A programming language used to write code.
      \item A formula for estimating running time.
    \end{enumerate}
  \end{block}

  \begin{exampleblock}{Question 2}
    Which statement best describes an \textbf{algorithm}?
    \begin{enumerate}
      \item A finite and precise procedure for solving a problem.
      \item A container used to store values.
      \item A Python library.
      \item A database table.
    \end{enumerate}
  \end{exampleblock}

  % Answers: Q1-B, Q2-A
\end{frame}

\begin{frame}{Algorithms use data structures to speed up operations}
  \small

  \begin{columns}[T]
    \begin{column}{0.48\textwidth}
      \begin{block}{A common scenario}
        Given a collection of records, determine whether an element $x$ is present. If it is present, retrieve the information associated with it.
      \end{block}

      \begin{exampleblock}{Real-world examples}
        \begin{itemize}
          \item Find a student using a student ID.
          \item Find a customer using a phone number.
          \item Check whether a username is already registered.
          \item Find a product and retrieve its price.
        \end{itemize}
      \end{exampleblock}
    \end{column}

    \begin{column}{0.48\textwidth}
      \begin{block}{A basic approach}
        Examine the records one by one:
        \begin{enumerate}
          \item Compare the current record with $x$.
          \item If they match, return the associated information. Otherwise, continue with the next record.
          \item Report that $x$ is not found after reaching the end.
        \end{enumerate}
      \end{block}

      \begin{alertblock}{What if the query is repeated?}
        But what happens when the collection is large and the same type of search must be performed thousands or millions of times?
      \end{alertblock}
    \end{column}
  \end{columns}
\end{frame}

\begin{frame}{Algorithms use data structures to speed up operations}
  \small

  \begin{beamercolorbox}[
      wd=\textwidth,
      sep=0.7em,
      center,
      rounded=true
    ]{block body example}
    A suitable data structure organizes the records in advance so that repeated search and retrieval operations can be performed more efficiently.
  \end{beamercolorbox}
\end{frame}

\begin{frame}{Algorithms use data structures to speed up operations}
  \small

  \begin{block}{Example: searching among one million records}
    Suppose the collection contains $n = 1{,}000{,}000$ records and the requested key is not present.
  \end{block}

  \begin{center}
    \begin{tabular}{p{0.27\textwidth}p{0.20\textwidth}p{0.40\textwidth}}
      \toprule
      \textbf{Organization} & \textbf{Approximate work} & \textbf{Explanation} \\
      \midrule
      Unsorted sequence & Up to $1{,}000{,}000$ comparisons & The algorithm may need to examine every record before concluding that the key is absent. \\
      Balanced binary search tree & About $20$ comparisons & Each comparison eliminates approximately half of the remaining search space: $\log_2(1{,}000{,}000) \approx 20$. \\
      \addlinespace
      Hash table & Often only a few probes on average & In a well-designed table, a search may require only about $2$--$3$ probes on average. \\
      \bottomrule
    \end{tabular}
  \end{center}
\end{frame}

\begin{frame}{Quick check: choosing a data structure}
  \small

  \begin{block}{Scenario}
    We need to store one million student records and repeatedly search by student ID.
  \end{block}

  \begin{alertblock}{Question}
    Which structure is usually more suitable?
    \begin{enumerate}
      \item Unsorted list
      \item Stack
      \item Hash table / dictionary
      \item Queue
    \end{enumerate}
  \end{alertblock}

  \begin{exampleblock}{Explain}
    Why can a suitable data structure reduce the number of comparisons?
  \end{exampleblock}

  % Answer: C
\end{frame}

\begin{frame}{Example of Data Structure 1}
  \begin{columns}[T]
    \begin{column}{0.49\textwidth}
      \begin{figure}
        \centering
        \includegraphics[width=1\linewidth]{figures/pixels_matrix.png}
        \caption{Image representation as a matrix}
      \end{figure}
    \end{column}
    \begin{column}{0.49\textwidth}
      \begin{figure}
        \centering
        \includegraphics[width=0.96\linewidth]{figures/user-rating.png}
        \caption{User-rating preferences matrix}
      \end{figure}
    \end{column}
  \end{columns}
\end{frame}

\begin{frame}{Example of Data Structure 2}
  \begin{columns}[T]
    \begin{column}{0.49\textwidth}
      \begin{figure}
        \centering
        \includegraphics[width=1\linewidth]{figures/trie.png}
        \caption{Words stored in a tree-based data structure}
      \end{figure}
    \end{column}
    \begin{column}{0.49\textwidth}
      \begin{figure}
        \centering
        \includegraphics[width=0.96\linewidth]{figures/dsu.png}
        \caption{Objects stored in an array-based but tree-like data structure}
      \end{figure}
    \end{column}
  \end{columns}
\end{frame}

\begin{frame}{Common operations and suitable data structures}
  \scriptsize

  \begin{center}
    \renewcommand{\arraystretch}{1.35}
    \begin{tabular}{p{0.25\textwidth}p{0.32\textwidth}p{0.31\textwidth}}
      \toprule
      \textbf{Common operation} & \textbf{Typical question} & \textbf{Suitable data structures} \\
      \midrule
      Store and access by position & What is the value at position $i$? & Array, list, matrix \\
      Search and membership & Does $x$ exist? What information is associated with key $x$? & Hash table, set, binary search tree \\
      Insert and remove elements & Which element should be added or removed next? & Linked list, stack, queue \\
      Select or prioritize elements & Which element is largest, smallest, or most urgent? & Heap, priority queue, binary search tree \\
      Represent hierarchical relationships & What is the parent, child, or descendant of an object? & Tree, binary tree \\
      Represent general relationships & Are two objects connected? Is there a path between them? & Graph \\
      \bottomrule
    \end{tabular}
  \end{center}
\end{frame}

\begin{frame}{Quick check: match operations to data structures}
  \scriptsize

  \begin{center}
    \renewcommand{\arraystretch}{1.35}
    \begin{tabular}{p{0.45\textwidth}p{0.38\textwidth}}
      \toprule
      \textbf{Operation} & \textbf{Suitable structure} \\
      \midrule
      Check whether a username already exists & \underline{\hspace{3.2cm}} \\
      Process customers in order of arrival & \underline{\hspace{3.2cm}} \\
      Select the most urgent task & \underline{\hspace{3.2cm}} \\
      Find words starting with a prefix & \underline{\hspace{3.2cm}} \\
      Represent roads between cities & \underline{\hspace{3.2cm}} \\
      \bottomrule
    \end{tabular}
  \end{center}

  % Suggested answers: set/hash table; queue; heap/priority queue; trie; graph
\end{frame}

\begin{frame}{Advanced operations and suitable data structures}
  \scriptsize

  \begin{center}
    \renewcommand{\arraystretch}{1.35}
    \begin{tabular}{p{0.25\textwidth}p{0.32\textwidth}p{0.31\textwidth}}
      \toprule
      \textbf{Common operation} & \textbf{Typical question} & \textbf{Suitable data structures / techniques} \\
      \midrule
      Merge and compare groups & Are $x$ and $y$ in the same group? How can two groups be joined? & Disjoint Set Union (DSU) \\
      Find common ancestors & What is the closest common ancestor of nodes $u$ and $v$? & LCA, binary lifting, Euler tour \\
      Query and update a range & What is the sum, minimum, or maximum from position $l$ to $r$? & Segment tree, Fenwick tree \\
      Calculate prefix information & What is the total from the beginning to position $i$? & Fenwick tree, prefix-sum array \\
      Search by prefix & Which words start with a given group of letters? & Trie \\
      Answer many queries on fixed data & What is the minimum or maximum value between positions $l$ and $r$? & Sparse table \\
      \bottomrule
    \end{tabular}
  \end{center}

  \vspace{0.3cm}

  \begin{alertblock}{Key idea}
    Advanced data structures are designed for special operations that appear many times in large or complex problems.
  \end{alertblock}
\end{frame}

\begin{frame}{Do we need to implement every data structure ourselves?}
  \small

  \begin{exampleblock}{Good news}
    Most programming languages already provide efficient implementations of many fundamental data structures.

    \vspace{0.35cm}

    In Python, for example:
    \begin{itemize}
      \item \texttt{list}: dynamic array and basic stack
      \item \texttt{dict}: hash table for key--value data
      \item \texttt{set}: hash-based collection of unique values
      \item \texttt{collections.deque}: queue and double-ended queue
      \item \texttt{heapq}: priority queue implemented using a heap
    \end{itemize}
  \end{exampleblock}

  \begin{beamercolorbox}[
      wd=\textwidth,
      sep=0.7em,
      center,
      rounded=true
    ]{block body example}
    We can use these data structures without implementing all their internal details from scratch.
  \end{beamercolorbox}
\end{frame}

\begin{frame}{Using a data structure still requires understanding}
  \small

  \begin{alertblock}{Bad news}
    Having data structures available does not automatically tell us:
    \begin{itemize}
      \item which data structure is suitable for a problem;
      \item which operations are efficient or inefficient;
      \item how much time and memory each operation requires;
      \item how the structure behaves in special cases.
    \end{itemize}
  \end{alertblock}

  \begin{block}{Example}
    Both a Python \texttt{list} and a \texttt{set} can store a collection of values. However:
    \begin{itemize}
      \item searching a list may require examining many elements;
      \item searching a set is usually much faster on average;
      \item a list preserves order, while a set does not.
    \end{itemize}
  \end{block}
\end{frame}

\begin{frame}{Using a data structure still requires understanding}
  \small

  \centering
  \textbf{\color{red}Choosing the wrong data structure can make an otherwise correct program inefficient.}
\end{frame}

\begin{frame}{Why do we study data structures?}
  \small

  \begin{columns}[T]
    \begin{column}{0.48\textwidth}
      \begin{block}{Use existing structures}
        We need to understand how to:
        \begin{itemize}
          \item select a suitable structure;
          \item use its operations correctly;
          \item analyze its time and memory costs;
          \item combine multiple structures.
        \end{itemize}
      \end{block}
    \end{column}

    \begin{column}{0.48\textwidth}
      \begin{exampleblock}{Develop new solutions}
        In some applications, existing structures may not directly support the required operations.

        \vspace{0.2cm}

        We may need to:
        \begin{itemize}
          \item adapt an existing structure;
          \item extend it with new operations;
          \item combine several structures;
          \item design a specialized data structure.
        \end{itemize}
      \end{exampleblock}
    \end{column}
  \end{columns}

  \vspace{0.4cm}

  \begin{beamercolorbox}[
      wd=\textwidth,
      sep=0.7em,
      center,
      rounded=true
    ]{block body example}
    \textbf{We study data structures not only to implement them, but also to select, use, analyze, adapt, and extend them appropriately.}
  \end{beamercolorbox}
\end{frame}

\begin{frame}{Quick check: Python built-in structures}
  \small

  \begin{block}{Choose a Python structure}
    \begin{enumerate}
      \item Store unique course codes and check membership quickly.
      \item Map student IDs to student names.
      \item Process tasks in first-in-first-out order.
      \item Repeatedly get the smallest priority value.
    \end{enumerate}
  \end{block}

  \begin{exampleblock}{Possible answers}
    \begin{enumerate}
      \item \texttt{set}
      \item \texttt{dict}
      \item \texttt{collections.deque}
      \item \texttt{heapq}
    \end{enumerate}
  \end{exampleblock}
\end{frame}

\begin{frame}{One problem, many possible algorithms}
  \small

  \begin{block}{Example problem in data science}
    Given a new data point $x$, find the $k$ most similar observations in a dataset.
  \end{block}

  \begin{exampleblock}{Possible algorithms}
    \begin{itemize}
      \item Linear scan: compare $x$ with every data point.
      \item KD-tree: organize low-dimensional data into searchable regions.
      \item Ball tree: organize points using distance-based regions.
      \item Approximate nearest-neighbor search: trade exactness for speed.
    \end{itemize}
  \end{exampleblock}

  \begin{alertblock}{Question}
    If several algorithms can solve the same problem, how do we choose the most suitable one?
  \end{alertblock}
\end{frame}

\begin{frame}{Choosing a suitable algorithm}
  \small

  \begin{columns}[T]
    \begin{column}{0.48\textwidth}
      \begin{block}{Input characteristics}
        \begin{itemize}
          \item Is the dataset small or large?
          \item Is the data low-dimensional or high-dimensional?
          \item Is exact search required?
          \item Are many queries repeated?
          \item Does the data change over time?
        \end{itemize}
      \end{block}
    \end{column}

    \begin{column}{0.48\textwidth}
      \begin{exampleblock}{Resource constraints}
        \begin{itemize}
          \item How much memory is available?
          \item How fast must each query be answered?
          \item Is preprocessing allowed?
          \item Is approximate output acceptable?
          \item Is implementation simplicity important?
        \end{itemize}
      \end{exampleblock}
    \end{column}
  \end{columns}

  \vspace{0.35cm}

  \begin{alertblock}{Key idea}
    Algorithm choice depends on the data, the constraints, and the required balance between accuracy, time, and memory.
  \end{alertblock}
\end{frame}

\begin{frame}{Why do we study algorithms and complexity?}
  \small

  \begin{block}{Algorithmic knowledge helps us answer}
    \begin{itemize}
      \item Which algorithm is correct for this problem?
      \item Which algorithm is fast enough for this input size?
      \item Which algorithm uses acceptable memory?
      \item Which algorithm behaves well in difficult cases?
      \item Which algorithm is suitable for the available data structure?
    \end{itemize}
  \end{block}

  \vspace{0.35cm}

  \begin{beamercolorbox}[
      wd=\textwidth,
      sep=0.75em,
      center,
      rounded=true
    ]{block body example}
    \textbf{A good algorithm is not only correct; it must also match the data, the constraints, and the performance requirements.}
  \end{beamercolorbox}
\end{frame}

\begin{frame}{Quick check: selecting an algorithm}
  \small

  \begin{block}{Scenario}
    In a recommendation system, given a new product, find the $k$ most similar products in a large dataset.
  \end{block}

  \begin{alertblock}{Questions}
    \begin{enumerate}
      \item Is this a search, sorting, graph, or nearest-neighbor problem?
      \item If the dataset is small, what simple approach can we use?
      \item If the dataset is very large, what constraints should we consider?
    \end{enumerate}
  \end{alertblock}

  \begin{exampleblock}{Key point}
    Algorithm choice depends on data size, dimensionality, memory, accuracy, and response-time requirements.
  \end{exampleblock}
\end{frame}

\begin{frame}{Data Structures, Algorithms and Maths help}
  \begin{columns}[T]
    \begin{column}{0.5\textwidth}
      \begin{figure}
        \centering
        \includegraphics[width=0.87\linewidth]{figures/DeepSeek-AI-vs-ChatGPT.jpg}
      \end{figure}
    \end{column}
    \begin{column}{0.5\textwidth}
      \begin{figure}
        \centering
        \includegraphics[width=0.9\linewidth]{figures/MIT-Optimal-Decisions.jpg}
      \end{figure}
      \centering{\url{https://arxiv.org/pdf/2505.21692}}
    \end{column}
  \end{columns}

  \vspace{0.2cm}

  \begin{beamercolorbox}[
      wd=\textwidth,
      sep=0.65em,
      center,
      rounded=true
    ]{block body example}
    \textbf{DSA and mathematics help us understand, verify, and improve AI-generated solutions instead of using them blindly.}
  \end{beamercolorbox}
\end{frame}

\begin{frame}{Course at a glance}
  \begin{columns}[T]
    \begin{column}{0.7\textwidth}
      \begin{block}{General information}
        \begin{itemize}
          \item \textbf{Title:} Data Structures and Algorithms in Python
          \item \textbf{Type:} Compulsory
          \item \textbf{Credits:} 3
        \end{itemize}
      \end{block}
    \end{column}
    \begin{column}{0.3\textwidth}
      \begin{exampleblock}{Study hours}
        \begin{itemize}
          \item Lectures: 45 hours
          \item Tutorials: 22.5 hours
          \item Self-study: 90 hours
        \end{itemize}
      \end{exampleblock}
    \end{column}
  \end{columns}

  \vspace{0.5cm}
  {\color{red} \bf Note:}
  \begin{itemize}
    \item Slides will be uploaded to the LMS before the class.
    \item There will be a tutorial every two lectures, scheduled by instructors.
    \item Office hours will be scheduled by instructors.
  \end{itemize}
\end{frame}

\begin{frame}{Instructors and learning resources}
  \begin{columns}[T]
    \begin{column}{0.45\textwidth}
      \begin{block}{Instructors}
        \small
        \textbf{Assoc. Prof. Dr. Nguyen Trung Thanh}\\
        \texttt{nguyenthanh@neu.edu.vn}

        \medskip
        \textbf{Dr. Vu Duc Minh}\\
        \texttt{minhvd@neu.edu.vn}

        \medskip
        \textbf{M.Sc. Nguyen Thanh Hoang}\\
        \texttt{hoangnt@neu.edu.vn}

        \medskip
        \textbf{M.Sc. Dam Tien Thanh}\\
        \texttt{thanhdt@neu.edu.vn}

        \medskip
        \textbf{M.Sc. Ta Dinh Quy}\\
        \texttt{quytd@neu.edu.vn}
      \end{block}
    \end{column}
    \begin{column}{0.55\textwidth}
      \begin{exampleblock}{Main textbook}
        \small
        Goodrich et al. (2013), \emph{Data Structures and Algorithms in Python}.
      \end{exampleblock}
      \begin{alertblock}{Additional resources}
        \small
        Cormen et al. (2022), \emph{Introduction to Algorithms}.

        \medskip
        Kleinberg and Tardos (2005), \emph{Algorithm Design}.

        \medskip
        Karumanchi (2015), \emph{Data Structure and Algorithmic Thinking with Python}.
      \end{alertblock}
    \end{column}
  \end{columns}
\end{frame}

\begin{frame}{What you learn}
  \small
  \begin{columns}[T]
    \begin{column}{0.48\textwidth}
      \begin{exampleblock}{Data structures}
        \begin{enumerate}
          \item Arrays and linked lists
          \item Stacks and queues
          \item Hash tables and heaps
          \item Trees and graphs
        \end{enumerate}
      \end{exampleblock}
    \end{column}
    \begin{column}{0.48\textwidth}
      \begin{block}{Algorithms}
        \begin{enumerate}
          \item Sorting and searching algorithms
          \item Traversal algorithms: DFS, BFS
          \item Shortest path and dynamic programming
          \item Other algorithms, time permitting
        \end{enumerate}
      \end{block}
    \end{column}
  \end{columns}

  \vspace{1cm}

  \centering{\color{red} Applications will be discussed primarily in the context of machine learning and data science.}
\end{frame}

\begin{frame}{Data structures --- An illustration}
  \begin{figure}
    \centering
    \includegraphics[width=1\linewidth]{figures/type_DS_fig.png}
    \caption{Examples of data structures studied in the course}
  \end{figure}
\end{frame}

\begin{frame}{Grading methods}
  \begin{center}
    \begin{tabular}{p{0.3\textwidth}p{0.4\textwidth}c}
      \toprule
      \textbf{Assessment} & \textbf{Format/Evaluation} & \textbf{Weight} \\
      \midrule
      Attendance & Based on class participation & \textbf{10\%} \\
      \addlinespace
      Homework assignment & Written paper & \textbf{20\%} \\
      \addlinespace
      Midterm exam & 01 written exam & \textbf{20\%} \\
      \addlinespace
      Final exam & 01 written exam & \textbf{50\%} \\
      \bottomrule
    \end{tabular}
  \end{center}

  {\color{red} \bf Note:}
  \begin{enumerate}
    \item This is not a programming course.
    \item There will likely be only a limited number of computer lab sessions.
  \end{enumerate}
\end{frame}

\begin{frame}{Attendance and exam formats}
  \begin{columns}[T]
    \begin{column}{0.55\textwidth}
      \begin{block}{Attendance}
        \begin{itemize}
          \item At least 80\% class attendance.
          \item Active participation in class.
        \end{itemize}
      \end{block}

      \begin{exampleblock}{Homework assignment}
        \begin{itemize}
          \item Every week.
          \item Handwritten on paper.
          \item Scan and upload to the LMS by the deadline.
        \end{itemize}
      \end{exampleblock}
    \end{column}

    \begin{column}{0.45\textwidth}
      \begin{exampleblock}{Midterm exam}
        \begin{itemize}
          \item Written exam: 90 minutes.
          \item Week 8, subject to change.
        \end{itemize}
      \end{exampleblock}
      \begin{alertblock}{Final exam}
        \begin{itemize}
          \item Written exam: 90 minutes.
          \item Week 17, subject to change.
          \item Condition: \textbf{Attendance\_score $\ge 5$}.
        \end{itemize}
      \end{alertblock}
    \end{column}
  \end{columns}
\end{frame}

\begin{frame}{How is DSA related to other courses?}
  \begin{tabular}{lp{8cm}}
    \toprule
    \textbf{Course} & \textbf{Main DSA Concepts} \\
    \midrule
    Algorithm Design & Dynamic programming, graphs, complexity \\
    Introduction to Database & B-trees, hashing, indexing \\
    Advanced Database & Query optimization, graph processing \\
    Introduction to AI & Trees, graphs, A*, BFS, DFS \\
    Information Security & Hashing, string matching \\
    Statistical Learning & KD-trees, heaps, optimization \\
    Data Mining & Trees, graphs, hash tables \\
    Big Data Technology & Distributed data structures, MapReduce \\
    Software Engineering & ADTs, complexity analysis \\
    Cloud Computing & Consistent hashing, distributed indexing \\
    \bottomrule
  \end{tabular}
\end{frame}

\begin{frame}{Key takeaway}
  \fcolorbox{black}{yellow!20}{%
    \begin{minipage}{\dimexpr\linewidth-2\fboxsep-2\fboxrule}
      {\bf If you skip DSA and go straight to AI frameworks, you may build applications faster---but without the fundamentals, you will struggle to reason about performance, debug failures, and design scalable solutions. DSA gives you the foundation to understand, evaluate, and improve AI-generated code.}
    \end{minipage}%
  }
\end{frame}

\begin{frame}{}
  {\Huge
  \centering{\bf Any Question?}
  }
\end{frame}

\end{document}
